Gain planning distributes gain, noise, and linearity along the receive chain so
that the system meets both its sensitivity and its dynamic-range goals.

\section{The Balancing Act}
Put too little gain early and later stages' noise dominates (poor sensitivity);
put too much and a strong signal compresses the front end (poor large-signal
handling). The plan sets each stage's gain, NF, and IP3 so the cascaded noise
figure and cascaded IP3 both land where they must.

\section{Sensitivity}
The minimum detectable signal is the noise floor plus the required SNR:
\[ \text{MDS} = -174\ \text{dBm/Hz} + 10\log_{10}B + \text{NF} + \text{SNR}_{min}. \]
Lowering NF (a good LNA first) or narrowing $B$ improves sensitivity.

\section{Spurious-Free Dynamic Range}
The top of the usable range is set by third-order intermodulation. The
spurious-free dynamic range --- the span from the noise floor to where IM3
products rise out of it --- is
\[ \text{SFDR} = \tfrac{2}{3}\,(\text{IIP3} - \text{MDS}). \]
Automatic gain control extends the \emph{instantaneous} range by backing off gain
on strong signals, but SFDR is the fundamental figure.
